\chapter{Conclusões}
\label{chap:conclusoes}

A pesquisa mostra que aproximação probabilística, compressão em frequência e substituição da resposta dispersiva afetam aspectos distintos da inferência da massa do gráviton. A comparação entre coeficientes Fourier, estimadores quadráticos e controles gaussianos permite identificar essas contribuições em experimentos pareados.

O primeiro resultado central é a insuficiência dos dois primeiros momentos para descrever, por uma normal, a distribuição dos estimadores físicos nos cenários estudados. A campanha tensorial apresentou dez rejeições de calibração mantidas sob a incerteza numérica considerada. Isso torna a adequação da verossimilhança uma questão física e estatística anterior à escolha de uma frequência de referência. Nos controles gaussianos corretamente especificados, a compressão reduziu a informação média sobre a massa em $0,06211\pm0,00477$ e $0,10631\pm0,00695$ nat nas duas configurações da população ampliada.

O segundo resultado é a distinção entre um limite superior e uma medida informativa. Com priori uniforme em $u\in[0,1]$ e verossimilhança constante, $Q_{0,95}(u)=0,95$ e $D_{\rm KL}=0$. O limite correspondente, $m_gc^2\leq0,95h_{\mathrm P}/T$, provém inteiramente do suporte. Comparar priori e posterior, quantificar o ganho de informação e verificar a sensibilidade à priori são, portanto, componentes essenciais da interpretação de limites de massa em PTAs.

A reprodução do TG estabeleceu as convenções de curvatura, unidades e vínculos que conectam a propagação massiva à resposta de temporização. A extensão de Fierz--Pauli preservou uma única amplitude de helicidade zero com deformações transversal e longitudinal vinculadas. No limiar, a proporcionalidade $\Gamma_0=\Gamma_T/2$ produz uma degenerescência angular entre as potências escalar e tensorial; sua separação depende de informação além da forma angular nesse limite, como a evolução espectral fora dele.

Os contrastes de frequência com marginalização dos parâmetros auxiliares permaneceram inconclusivos na escala de 0,01 da largura da priori. A redução a uma inferência condicional de massa, em 50 realizações com ruído fixado, permitiu resolver a sensibilidade numérica e encontrar pequenos deslocamentos médios de $Q_{95}(u)$: os quatro intervalos das médias ficaram dentro de $[-0{,}01;0{,}01]$. A comparação identifica uma situação em que a aproximação é pouco influente nessa estatística e concentra a questão restante no papel da marginalização.

O alcance atual é o de simulações com geometria controlada e uma extensão com posições reais de catálogo. O experimento de ajuste de timing mostrou que a projeção dos parâmetros de rotação e dos termos anuais induz correlações entre frequências e pseudocovariância, mesmo com ruído gaussiano. Sua incorporação à inferência conjunta de massa e ruído constitui o passo seguinte para aplicações observacionais. Os resultados fornecem critérios concretos para essa passagem: testar a distribuição dos estimadores, separar perda de informação de erro de resposta e distinguir restrições impostas pelos dados das induzidas pela priori.
